\documentclass[11pt]{article}
\usepackage[margin=2.5cm]{geometry}
\usepackage{amsmath}
\newcommand{\kev}{\,\mathrm{keV}}
\title{spexai: analysis decision log}
\date{}
\begin{document}
\maketitle

Appended as choices are made: what was tried, what was chosen, and why.

\section*{2026-09-15: DEM log-$T$ switch (bias campaign)}

\paragraph{Parametrisation.} \emph{Chosen:} Gaussian in $\log_{10}T$
(SPEX \texttt{gdem}) for the bias campaign. \emph{Tried:} Gaussian in linear
$T$ (thesis parametrisation). \emph{Why:} on the log-spaced grid a log-$T$
width has constant resolution; the linear-$T$ width spanned 0.6--20 cells across
one design and produced the single $>1\sigma$ outlier (1.642) as a resolution
artefact.

\paragraph{Temperature design range.} \emph{Chosen:} $0.7$--$15\kev$ for both
the DEM centre and single-$T$ $kT$, drawn uniform in $\log_{10}T$.
\emph{Tried/considered:} $1.5$--$8\kev$ (previous, uncited), $1$--$12\kev$,
$0.5$--$15\kev$. \emph{Why:} groups through the most massive clusters.
$0.5\kev$ rejected: below the PCHIP truth's safe floor ($0.7\kev$; Ar diverges
at $0.5\kev$), below the emulator floor ($0.5013\kev$), and where the cache step
defects cluster (Na, Fe, Ni, Cu, Co at $0.52$--$0.67\kev$).

\paragraph{Width range.} \emph{Chosen:} $\sigma_x = 0.056$--$0.4$\,dex, no cap
tied to the centre. \emph{Tried:} a centre-dependent cap requiring $\ge98\%$ of
the DEM on the grid. \emph{Why rejected:} SPEX does not restrict \texttt{gdem}
to fully contained Gaussians, and truth and emulator share the same
$0.7$--$19.94\kev$ grid, so a truncated DEM is truncated identically on both
sides and $b_{\rm sys}$ remains a fair comparison. The floor is a resolution
limit ($\approx2.7$ cells of the 70-node grid). The contained fraction is
recorded per truth point for interpretation instead.

\paragraph{Fit bounds and steps.} \emph{Chosen:} $\log T$ centre and $\log kT$:
design range $\pm0.1$\,dex; $\sigma_x$: $0.045$--$0.5$\,dex; finite-difference
step $5\times10^{-4}$\,dex for both DEM parameters. \emph{Why:} bounds are not
physics; they only need never to bind at a truth point and to stay inside the
emulator's trained range. The width's lower bound sits just below the
resolution floor. Single-$T$ $kT$ keeps its $5\times10^{-3}\kev$ step
(unchanged; to revisit if the step check flags it at $0.7\kev$).

\paragraph{Campaign fiducial.} \emph{Chosen:} $\log_{10}T = 0.591$,
$\sigma_x = 0.111$\,dex, converted from $3.9\pm1.0\kev$ at the centre
($\sigma_x = \sigma_T/(T\ln 10)$).

\paragraph{Perseus showcase.} \emph{Chosen:} stays linear $T$. \emph{Why:} the
XRISM Perseus enrichment analysis (arXiv:2606.17141) fits XSPEC
\texttt{bvvgadem}, whose $\sigma_T$ is in keV; the showcase should match the
real analysis. Grid and fiducial deferred.

\paragraph{Frozen scripts.} \emph{Chosen:} hot\_floor scripts, \texttt{dump\_truth},
\texttt{bake\_off} and the cache-step diagnostics are not modified; their
divergence from the campaign is recorded in \texttt{NOTES.md}. \emph{Why:} not
needed for P6/P7, and their produced results must remain reproducible.
Consequence: \texttt{campaign.gaussian\_dem} and \texttt{campaign.build\_params}
keep the linear-$T$ path; the campaign uses new \texttt{gaussian\_logT\_dem}.

\paragraph{Finite-difference step (2026-09-16).} \emph{Chosen:} one step in
dex, $5\times10^{-4}$, for every temperature axis; $kT$ takes it converted at
the point, $kT\ln 10\,\times 5\times10^{-4}$, replacing the inherited absolute
$5\times10^{-3}\kev$. \emph{Why:} the step check at the widened range showed the
DEM parameters comfortably in the valley between truncation ($\propto h^2$) and
float32 round-off ($\propto1/h$) --- halving the step moved the Fisher-weighted
Jacobian column by $7.7\times10^{-5}$ to $8.3\times10^{-4}$ --- while single-$T$
$kT$ moved by $6.0\times10^{-2}$ at $0.75\kev$ against $3.1\times10^{-3}$ at
$14\kev$. An absolute step is $0.67\%$ of $kT$ at the cold end and $0.036\%$ at
the hot end, and above $1.9\kev$ a cold plasma contributes only its exponential
tail, so the curvature there is extreme; the implied Jacobian bias at the old
step is ${\sim}8\%$. \texttt{campaign.build\_params} keeps $5\times10^{-3}\kev$
(frozen hot\_floor path, fiducial $3.9\kev$, where it is $0.13\%$ of $kT$).

\paragraph{Outcome of that change (2026-09-16), and a correction.} The rerun
refuted the reasoning above. With the relative step, single-$T$ $kT$ at
$0.75\kev$ moved from $6.8\times10^{-2}$ to $3.8\times10^{-2}$ and at $14\kev$
from $3.1\times10^{-3}$ to $6.6\times10^{-3}$; the DEM values were unchanged.
Shrinking the step $5.8\times$ at the cold end should have reduced the
difference ${\sim}34\times$ under an $h^2$ truncation law and reduced it
$1.8\times$, and neither end follows the round-off law either. Step size is
therefore not what controls the cold-end number, and the ${\sim}8\%$ Jacobian
bias inferred above does not follow from the data. \emph{The relative step is
retained} on the grounds of units---a single step in dex serves $0.7$--$15\kev$
where an absolute keV step cannot---not on that measurement.

What remains unexplained is confined to the cold points. Above $1.9\kev$ a
$0.75\kev$ plasma leaves nearly every channel with $\mu\simeq0$, and both the
check and \texttt{linear\_bias\_fisher} weight channels by ${\sim}1/\mu$, so
empty channels can dominate on float32 noise alone. Whether $J$ is genuinely
unreliable below ${\sim}1.5\kev$, the check misreports it, or the band simply
carries no information there is not resolved; the diagnostic that separates
them (three step sizes, a signal-restricted metric, and a count of empty
channels) is specified but was not run. This affects only the coldest Tier~B
points, where $\mathrm{cond}(F)$ is reported per point; $k$ is insulated,
because the Gauss--Newton fixed point is set by the autograd score and $F$ only
preconditions the path.

\paragraph{Forward-mode AD instead of a step: probed, rejected for now.}
\emph{Tried:} replacing the central-difference Jacobian with
\texttt{torch.func.jvp}. Forward mode is the right shape for this object (one
pass per parameter, versus one per channel for reverse mode), it would be about
twice as cheap as the $2n+1$ stencil, and it has no step at all.
\emph{Result:} it fails on both paths --- \texttt{Sparse CSR tensors do not have
strides} in the response fold (DEM), and \texttt{You must implement the jvp
function for custom autograd.Function} from the trunk's gradient checkpointing
(single-$T$). An abundance control agreed to $5.8\times10^{-4}$ by finite
differences in both modes, so the probe itself was sound.
\emph{Why not fixed now:} both are fixable --- the fold is linear, so forward
mode need only run as far as \texttt{flux} and the same map applies to the
tangent, and checkpointing is redundant under forward mode --- but it is a
hot-path restructure with an uncertain interaction with \texttt{grad\_enabled},
to fix one parameter at one end of the design that the relative step already
fixes. Recorded as future work; probe kept at
\texttt{scripts/inference/probe\_jacobian\_ad.py}.

\section*{2026-09-16: Gauss--Newton bugs (D1--D3)}

\paragraph{D1, vacuous convergence: FIXED.} \emph{Chosen:} \texttt{converged}
now requires \texttt{resolved.any()}, and the record carries
\texttt{n\_resolved}. \emph{Why:} \texttt{frac} is zeroed wherever a
parameter's bias is unresolved, so a point that resolved nothing scored
$\max(\texttt{frac})=0$ and reported converged---``the fit learned nothing''
became ``the fit converged''. All six converged DEM points of the 2026-09-10
run were that case, each clamp-saturated at $\texttt{drift\_sigma}=20.000$, and
their $k$ was read as a result. The criterion was extracted from
\texttt{run\_point} into \texttt{p6\_sweep.convergence\_verdict} so it could be
tested at all; three tests added.

\paragraph{D2, Gauss--Newton does not descend: HELD.} $-\log L$ rises
monotonically while the Newton decrement falls. \emph{Deferred by D.H.} until
the log-$T$ switch and D1 have been through a real screen, on the grounds that
the other fixes may change the picture. The proposed fix, if it is still needed,
is to accept a step only on decrease and otherwise halve along the same Newton
direction---affine invariant, and not the strong-Wolfe search that broke
L-BFGS.

\paragraph{D3, unreachable spread criterion: more iterations.} \emph{Chosen:}
raise the iteration budget (\texttt{--gn\_iter} $8\to20$), plus a start-up
check that warns when $\texttt{gn\_iter}\times\texttt{gn\_max\_step}$ is smaller
than the point's actual start spread. \emph{Why:} the noiseless starts scatter
by ${\sim}|b_{\rm sys}|$ per component---$60$--$280\sigma$ between the most
distant pair---while the trust region caps each iteration at $20\sigma$, so
closing the gap needs $\ge14$ iterations before convergence can begin and the
budget was $8$--$12$. The criterion therefore failed by construction. This is
affordable only now: at ${\sim}12$\,s per iteration $20$ iterations is
${\sim}4$ minutes per point, where at $980$\,s it was $3.3$ hours.
\emph{Rejected:} a larger step clamp (trades a known cost for overshoot risk
into a region where the dropped residual term is not small) and tighter starts
via \texttt{--start\_bsys\_scale} (closer starts agree more easily, so the
multi-start certificate would certify less). Note the budget only helps if the
iteration descends, i.e.\ D2; both should be re-read against the new screen.

\paragraph{TwoGaussianDEM.} \emph{Chosen:} no longer renormalised, matching
\texttt{ParametricDEM}. \emph{Why:} renormalising silently redistributes an
off-grid component's emission measure onto the grid.

\end{document}
